% Options for packages loaded elsewhere
\PassOptionsToPackage{unicode}{hyperref}
\PassOptionsToPackage{hyphens}{url}
\documentclass[
  11pt,
]{article}
\usepackage{xcolor}
\usepackage[margin=1in]{geometry}
\usepackage{amsmath,amssymb}
\setcounter{secnumdepth}{-\maxdimen} % remove section numbering
\usepackage{iftex}
\ifPDFTeX
  \usepackage[T1]{fontenc}
  \usepackage[utf8]{inputenc}
  \usepackage{textcomp} % provide euro and other symbols
\else % if luatex or xetex
  \usepackage{unicode-math} % this also loads fontspec
  \defaultfontfeatures{Scale=MatchLowercase}
  \defaultfontfeatures[\rmfamily]{Ligatures=TeX,Scale=1}
\fi
\usepackage{lmodern}
\ifPDFTeX\else
  % xetex/luatex font selection
  \setmainfont[]{Times New Roman}
\fi
% Use upquote if available, for straight quotes in verbatim environments
\IfFileExists{upquote.sty}{\usepackage{upquote}}{}
\IfFileExists{microtype.sty}{% use microtype if available
  \usepackage[]{microtype}
  \UseMicrotypeSet[protrusion]{basicmath} % disable protrusion for tt fonts
}{}
\makeatletter
\@ifundefined{KOMAClassName}{% if non-KOMA class
  \IfFileExists{parskip.sty}{%
    \usepackage{parskip}
  }{% else
    \setlength{\parindent}{0pt}
    \setlength{\parskip}{6pt plus 2pt minus 1pt}}
}{% if KOMA class
  \KOMAoptions{parskip=half}}
\makeatother
\usepackage{longtable,booktabs,array}
\newcounter{none} % for unnumbered tables
\usepackage{calc} % for calculating minipage widths
% Correct order of tables after \paragraph or \subparagraph
\usepackage{etoolbox}
\makeatletter
\patchcmd\longtable{\par}{\if@noskipsec\mbox{}\fi\par}{}{}
\makeatother
% Allow footnotes in longtable head/foot
\IfFileExists{footnotehyper.sty}{\usepackage{footnotehyper}}{\usepackage{footnote}}
\makesavenoteenv{longtable}
\setlength{\emergencystretch}{3em} % prevent overfull lines
\providecommand{\tightlist}{%
  \setlength{\itemsep}{0pt}\setlength{\parskip}{0pt}}
\usepackage{bookmark}
\IfFileExists{xurl.sty}{\usepackage{xurl}}{} % add URL line breaks if available
\urlstyle{same}
\hypersetup{
  hidelinks,
  pdfcreator={LaTeX via pandoc}}

\author{}
\date{}

\begin{document}

\section{Response to the Academic
Editor}\label{response-to-the-academic-editor}

\textbf{Manuscript:} \emph{Proxy-Based Diagnostics of Quantum Oracle
Sketching Robustness for Non-IID Sensor and Telemetry Streams}\\
\textbf{Journal:} Sensors (MDPI) · \textbf{Manuscript ID:}
sensors-4470240

\begin{center}\rule{0.5\linewidth}{0.5pt}\end{center}

\textbf{Editor's comment.} \emph{``I strongly invite the Authors to not
include preprints among references. These works are still undergoing the
peer-review process, thus meaning that their content may change prior to
publication (if they will be finally published). The Authors can
substitute them with related works already published.''}

\emph{Response.} We thank the Academic Editor for this instruction. We
agree with the principle behind it and have acted on it: the entire
reference list was audited entry by entry, every preprint that has since
been published has been replaced by its peer-reviewed version, and the
reference list now contains \textbf{one} preprint, which is the specific
framework whose robustness this paper studies. Below we report what the
audit found, explain why that single entry cannot be substituted without
misattributing its results, and describe the safeguards --- including a
new footnote added in response to this comment --- that ensure no
conclusion of this paper depends on it.

\subsection{1. What the audit found}\label{what-the-audit-found}

Every one of the 54 references was classified by publication type
(journal article, conference paper, preprint, report, thesis, book,
website) and then verified at the metadata level: all 51 entries
carrying a DOI were resolved through doi.org content negotiation and
machine-compared against the cited title, first author, venue, and year;
the three entries without a DOI were checked against their publisher's
own record.

{\def\LTcaptype{none} % do not increment counter
\begin{longtable}[]{@{}ll@{}}
\toprule\noalign{}
Outcome & Count \\
\midrule\noalign{}
\endhead
\bottomrule\noalign{}
\endlastfoot
References audited & 54 \\
Preprints present before the audit & 2 \\
\textbf{Preprints replaced with the published version} & \textbf{1} \\
Preprints remaining (no published version exists) & 1 \\
DOIs resolved and metadata-matched & 51 / 51 \\
Metadata errors found and corrected & 1 \\
\end{longtable}
}

\textbf{Replaced.} Kallaugher, Parekh and Voronova, \emph{How to design
a quantum streaming algorithm without knowing anything about quantum
computing}, previously cited as an arXiv preprint, has since appeared in
the \emph{Proceedings of the 2025 Symposium on Simplicity in Algorithms
(SOSA)}, pp.~9--45. The entry now carries the publisher DOI
10.1137/1.9781611978315.2, and the author order was confirmed against
the publisher's article page and the proceedings table of contents.

\textbf{Corrected.} The audit also caught an unrelated defect that would
otherwise have reached print: the entry for Su, Guo and Wang (2024)
carried a transposed article number whose DOI resolved to a different
paper. It now reads \emph{Information Sciences} \textbf{676}, article
120799, DOI 10.1016/j.ins.2024.120799, verified against the publisher's
record.

\subsection{2. The one preprint that
remains}\label{the-one-preprint-that-remains}

The remaining entry is Zhao et al., \emph{Exponential quantum advantage
in processing massive classical data} (arXiv:2604.07639), labelled
``Preprint'' in the bibliography.

We confirmed --- and, in preparing this reply, re-confirmed --- that no
peer-reviewed version exists to substitute. As of 16 August 2026: the
arXiv record remains at version 1 with no journal reference and no
related DOI; Crossref, Semantic Scholar, dblp, INSPIRE-HEP, Google
Scholar, and OpenReview all record the work as arXiv-only; it appears in
none of the accepted-papers lists of FOCS 2026, STOC 2026, QIP 2026, TQC
2026, CCC 2026, or SODA 2026, and on no publisher page of
\emph{Quantum}, \emph{PRX Quantum}, \emph{npj Quantum Information},
\emph{Nature}, or \emph{Science}; the publication pages of all seven
authors list it as arXiv-only; and none of the works citing it records a
venue or an in-press status for it.

We respectfully submit that this entry is a different case from the one
the Editor's comment is aimed at. It is not cited as supporting evidence
for a claim of ours, where a related published work could stand in its
place. \textbf{It is the object of study.} This paper examines how the
correlation assumptions of that specific framework --- its refreshing
time and repetition number, its ×R sample overhead, its task-specific
classical hardness results --- behave when the data stream is not IID.
Substituting a related published work would attribute those particular
theorems to authors who did not prove them, and removing the citation
would leave the framework we analyse unattributed. Either course would
introduce a citation-integrity problem more serious than the one the
instruction is intended to prevent.

\subsection{3. Why no conclusion of this paper depends on
it}\label{why-no-conclusion-of-this-paper-depends-on-it}

The Editor's concern is that a preprint's content may change before
publication. We would like to show concretely why that risk does not
propagate into this paper's findings:

\begin{enumerate}
\def\labelenumi{\arabic{enumi}.}
\tightlist
\item
  \textbf{No quantum algorithm is implemented or simulated.} The
  manuscript states this verbatim in the Abstract, Section 3.3, four
  separate results subsections, the Threats to Validity section, and the
  Conclusions. Nothing in the paper executes, reproduces, or empirically
  confirms any result of the preprint.
\item
  \textbf{Every inherited quantity is tabulated with its provenance.}
  Table 1 (Section 3.3) lists, per quantity, what is inherited from the
  framework, what this paper introduces, and what is \emph{not}
  theoretically guaranteed. If any theorem of the preprint changes, the
  affected rows are identifiable at a glance rather than diffused
  through the text.
\item
  \textbf{The paper's own contribution is independent of it.} The (τ, r)
  sensitivity landscape, the classical streaming baselines, the
  correlation diagnostics, and the two real-telemetry case studies are
  classical measurements that stand on their own regardless of the
  preprint's fate.
\item
  \textbf{The revision already removed any dependence on the framework's
  claims being settled.} Following the reviewers' recommendations, the
  manuscript was repositioned as a proxy-based diagnostic and
  hypothesis-generation framework; it treats the framework's results as
  premises to be stress-tested, not as established facts to be relied
  upon, and states plainly that it ``demonstrates no quantum advantage
  and implements no quantum algorithm''.
\item
  \textbf{Its theorems are never invoked alone.} Each invocation is
  accompanied by peer-reviewed anchors: Kallaugher, \emph{A quantum
  advantage for a natural streaming problem} (IEEE FOCS 2021,
  pp.~897--908); Kallaugher, Parekh and Voronova (SOSA 2025, pp.~9--45);
  and Huang et al., \emph{Quantum advantage in learning from
  experiments} (\emph{Science} \textbf{376}, 1182--1186, 2022). These
  are precisely the ``related works already published'' the Editor
  suggests, and they now carry the general claims about quantum
  streaming separations wherever such claims are made.
\end{enumerate}

\subsection{4. What we have changed in response to this
comment}\label{what-we-have-changed-in-response-to-this-comment}

Beyond the audit, we have reduced the manuscript's reliance on the
preprint to the minimum that honest attribution permits, and
strengthened the published anchors around it:

\begin{enumerate}
\def\labelenumi{\arabic{enumi}.}
\tightlist
\item
  \textbf{Citation reduction (eight occurrences → six).} Two repeat
  citations carried no independent attribution and were consolidated: a
  second citation within the same paragraph of Section 3.3 now reads
  ``the originating framework'', and the quantum-memory bullet in
  Section 5.3 now points to Section 3.3, where the provenance is
  established. The preprint is now cited only where attributing its
  specific results requires it.
\item
  \textbf{A stronger published anchor added.} We added Kallaugher,
  Parekh and Voronova, \emph{Exponential quantum space advantage for
  approximating maximum directed cut in the streaming model}
  (Proceedings of the 56th Annual ACM Symposium on Theory of Computing,
  STOC 2024, pp.~1805--1815, DOI 10.1145/3618260.3649709) --- the first
  exponential quantum--classical space separation for a natural
  streaming problem, and the strongest peer-reviewed result of the kind
  the preprint extends. It is co-cited in the Related Work and in the
  disclosure footnote (bibliography now 55 entries;
  citation--bibliography consistency re-verified).
\item
  \textbf{A corrected attribution.} In verifying the anchors we found
  that our Related Work sentence credited Kallaugher's 2021 FOCS paper
  with the \emph{first exponential} separation; that paper's advantage
  is polynomial (the exponential result for a natural streaming problem
  is the STOC 2024 paper above). The sentence now attributes each result
  correctly --- a further accuracy improvement prompted directly by the
  Editor's comment.
\end{enumerate}

At the first citation of the preprint in the Introduction we have also
added a footnote making its status explicit to every reader:

\begin{quote}
``This reference is a preprint (arXiv:2604.07639) and, at the time of
writing, had not appeared in a peer-reviewed venue; it is labelled as
such in the bibliography. It is cited because it is the specific
framework whose robustness under correlated data this paper examines, so
no already-published work can stand in for it without misattributing its
results. No conclusion of this paper depends on its final published
form: no quantum algorithm is implemented or simulated here, every
quantity inherited from the framework is listed with its provenance and
epistemic status in Table 1, and its theorems are invoked alongside
peer-reviewed results on quantum streaming separations and on
experimentally demonstrated quantum advantage.''
\end{quote}

We also undertake to monitor the record and, should the work appear in a
peer-reviewed venue before this manuscript reaches proofs, to replace
the entry with the published version at that stage.

\subsection{5. If the Editor still prefers
removal}\label{if-the-editor-still-prefers-removal}

We recognise that this remains the Editor's decision. We would only note
that removing the citation entirely would require either dropping the
attribution for the framework the paper analyses, or presenting its
assumptions as our own --- neither of which we are willing to do. If the
Editor judges that the single labelled preprint is nevertheless
unacceptable, we would be grateful for guidance on the preferred remedy,
and we will follow it.

We are grateful for the attention the Editor has given the reference
list; the audit prompted by this comment improved the manuscript beyond
the preprint question itself, and every one of the 54 references has now
been verified against its publisher's record.

On behalf of all authors,

\textbf{AbdelMoniem Helmy} (corresponding author)
abdelmoniem.hafez@cu.edu.eg · ORCID 0000-0001-5996-6019

\end{document}
